% =====================================================================
%  2bat-ud2-15.tex  —  SESSIÓ 15: Taller de preparació de la prova
% =====================================================================
\horatitol{Sessió 15 — Taller de preparació de la prova}%
{Conèixer l'estructura de la prova i la rúbrica, assajar un simulacre per blocs i fer un pla d'estudi personal.}

\subsection{Idea clau}

Ens preparem per a la prova \textbf{sense sorpreses}. La prova de la sessió següent té
\textbf{quatre blocs}, un per cada criteri (C1-C4). Saber exactament què es demana i
com es corregeix fa que l'examen deixi de ser imprevisible. Després assajarem amb un
\textbf{simulacre}.

\begin{idea}
\textbf{Estructura de la prova} (un bloc per criteri):
\begin{enumerate}[leftmargin=1.6em, itemsep=2pt]
  \item \textbf{Bloc 1 (C1) — Traduir:} passar un enunciat social a un sistema
        d'inequacions amb dues variables.
  \item \textbf{Bloc 2 (C2) — Representar:} dibuixar la regió factible i trobar-ne els
        vèrtexs.
  \item \textbf{Bloc 3 (C3) — Optimitzar:} avaluar la funció objectiu als vèrtexs i
        triar el màxim o el mínim.
  \item \textbf{Bloc 4 (C4) — Interpretar:} llegir la solució òptima en el context i,
        si escau, reconèixer un cas especial.
\end{enumerate}
\textbf{Es valora}, a més del resultat: la \textbf{traducció acurada} de les
condicions, el \textbf{dibuix net} de la regió i la \textbf{interpretació} dins del
context.
\end{idea}

\subsection{Rúbrica simplificada}

\noindent Així es corregeix cada bloc (nivells d'assoliment):
\begin{center}
\renewcommand{\arraystretch}{1.3}
\begin{tabular}{p{2.4cm} p{10.2cm}}
\toprule
\textbf{Nivell} & \textbf{Què demostra} \\
\midrule
\textbf{AE} (excel·lent) & Resol sense errors, representa amb cura i interpreta el
resultat dins del context. \\
\textbf{AN} (notable) & Resol correctament amb algun error menor; representació i
interpretació essencialment bones. \\
\textbf{AS} (suficient) & Resol el procediment bàsic amb alguns errors; representació
o interpretació incompletes. \\
\textbf{NA} (no assolit) & No identifica el procediment o comet errors que
n'invaliden el resultat. \\
\bottomrule
\end{tabular}
\end{center}

\subsection{Simulacre — un ítem de cada bloc}

\begin{exemple}[com es resol un ítem ben fet]
\textbf{Ítem tipus (Bloc 3):} sobre els vèrtexs $(0,0)$, $(5,0)$, $(2,4)$, $(0,6)$,
maximitza $P=3x+2y$. Un \textbf{ítem complet} avalua \emph{tots} els vèrtexs i tria:
\[
P(0,0)=0,\ P(5,0)=15,\ P(2,4)=14,\ P(0,6)=12 \ \Rightarrow\ \text{màxim } P=15 \text{ a } (5,0).
\]
\emph{Avaluar tots els vèrtexs i dir clarament on és l'òptim distingeix un AE d'un AS.}
\end{exemple}

\begin{exercicis}
\textbf{Resol el simulacre amb temps controlat. Un ítem per bloc:}
\begin{exlist}
  \item \textbf{(Bloc 1 — C1)} Tradueix a un sistema: «una entitat fa $x$ tallers i $y$
        sortides; com a màxim $12$ activitats en total i com a màxim $5$ sortides; res
        negatiu».
  \item \textbf{(Bloc 2 — C2)} Representa la regió factible de
        $\;x+y\le 8,\;x+2y\le 12,\;x\ge 0,\;y\ge 0\;$ i dóna'n els vèrtexs.
  \item \textbf{(Bloc 3 — C3)} Sobre els vèrtexs de l'exercici anterior, maximitza
        $P=4x+5y$.
  \item \textbf{(Bloc 4 — C4)} Interpreta la solució òptima del bloc 3 en termes de
        l'enunciat (quantes activitats de cada tipus i quantes «unitats» de l'objectiu).
\end{exlist}
\end{exercicis}

\subsection{Formulari-resum (suport)}

\begin{idea}
\textbf{Tingues-ho a mà mentre prepares la prova:}
\begin{itemize}[leftmargin=1.4em, itemsep=2pt]
  \item \textbf{Inequació amb dues variables:} recta frontera (igualtat) + tipus de
        línia (contínua $\le,\ge$; discontínua $<,>$) + semiplà (punt de prova, sovint
        l'origen).
  \item \textbf{Regió factible:} intersecció de tots els semiplans; recorda
        $x\ge 0$, $y\ge 0$.
  \item \textbf{Vèrtex:} creuament de dues vores; es troba \textbf{resolent el
        sistema} de les dues rectes.
  \item \textbf{Òptim:} avalua la funció objectiu a \textbf{cada} vèrtex; tria el més
        gran (màxim) o el més petit (mínim).
\end{itemize}
\end{idea}

\noindent\textbf{Pla d'estudi personal.} A partir de la teva graella d'autoavaluació
(sessió 7) i del simulacre, anota els \textbf{dos o tres aspectes} que has de repassar
amb més atenció:
\begin{center}
\renewcommand{\arraystretch}{1.5}
\begin{tabular}{c p{10.5cm}}
\toprule
\textbf{\#} & \textbf{Què he de repassar (i com)} \\
\midrule
1 & \rule{9.8cm}{0.4pt} \\
2 & \rule{9.8cm}{0.4pt} \\
3 & \rule{9.8cm}{0.4pt} \\
\bottomrule
\end{tabular}
\end{center}

% ---------------------------------------------------------------------
\fullsolucions{Sessió 15 — simulacre}
\begin{solucions}
\begin{sollist}
  \item \textbf{(Bloc 1)} $x+y\le 12$, \;$y\le 5$, \;$x\ge 0$, \;$y\ge 0$.
  \item \textbf{(Bloc 2)} Vèrtexs: $(0,0)$; sobre $y=0$ mana $x\le 8$ $\rightarrow
        (8,0)$; creuament $x+y=8$ i $x+2y=12$ ($y=4$, $x=4$) $\rightarrow (4,4)$; sobre
        $x=0$ mana $y\le 6$ $\rightarrow (0,6)$.
  \item \textbf{(Bloc 3)} $P=4x+5y$: $P(0,0)=0$, $P(8,0)=32$, $P(4,4)=16+20=36$,
        $P(0,6)=30$. Màxim $P=36$ a $(4,4)$.
  \item \textbf{(Bloc 4)} L'òptim $(4,4)$ vol dir fer \textbf{$4$ tallers i $4$
        sortides}, que dóna $P=36$ unitats de l'objectiu (la màxima possible). En aquest
        punt s'esgoten tots dos recursos ($x+y=8$ i $x+2y=12$).
\end{sollist}
\end{solucions}
